\textbf{Phần 1. Trắc nghiệm lựa chọn}
\setcounterpageref{ex}{0}

\begin{ex} %[4D8N46]
Muốn tìm số trung bình cộng của nhiều số, ta
\choice
{\True Tính tổng của các số đó rồi chia tổng đó cho số các số hạng}
{Lấy tổng của các số đó nhân với số các số hạng}
{Lấy số lớn nhất cộng số bé nhất chia đôi}
{Lấy hiệu của các số chia cho số các số hạng}
\loigiai{
Trung bình cộng $=$ Tổng các số $:$ Số các số hạng.\\
\textbf{Đáp án của câu hỏi là A}
}
\end{ex}

\begin{ex} %[4D8N46]
Trung bình cộng của hai số $40$ và $60$ là
\choice
{$45$}
{$55$}
{\True $50$}
{$100$}
\loigiai{
$(40 + 60) : 2 = 50$.\\
\textbf{Đáp án của câu hỏi là C}
}
\end{ex}

\begin{ex} %[4D8H46]
Số trung bình cộng của 3 số: $15; 25; 35$ là
\choice
{$20$}
{\True $25$}
{$30$}
{$75$}
\loigiai{
$(15 + 25 + 35) : 3 = 25$.\\
\textbf{Đáp án của câu hỏi là B}
}
\end{ex}

\begin{ex} %[4D8H46]
Bốn bạn có số bi lần lượt là $12, 14, 16, 18$. Số bi trung bình của mỗi bạn là
\choice
{$14$ viên}
{$16$ viên}
{$17$ viên}
{\True $15$ viên}
\loigiai{
$(12 + 14 + 16 + 18) : 4 = 60 : 4 = 15$ viên.\\
\textbf{Đáp án của câu hỏi là D}
}
\end{ex}

\begin{ex} %[4D8H46]
Trung bình cộng của hai số là $45$. Tổng của hai số đó là
\choice
{$45$}
{\True $90$}
{$135$}
{$22{,}5$}
\loigiai{
Tổng $= 45 \times 2 = 90$.\\
\textbf{Đáp án của câu hỏi là B}
}
\end{ex}

\begin{ex} %[4D8H46]
Trung bình cộng của 5 số lẻ liên tiếp là 15. Số chính giữa trong dãy 5 số đó là
\choice
{$13$}
{$17$}
{\True $15$}
{$11$}
\loigiai{
Trung bình cộng của dãy số cách đều lẻ số phần tử chính là số ở chính giữa.\\
\textbf{Đáp án của câu hỏi là C}
}
\end{ex}

\begin{ex} %[4D8V46]
Một ô tô giờ thứ nhất chạy được $45\text{ km}$, giờ thứ hai chạy được $55\text{ km}$, giờ thứ ba chạy được $50\text{ km}$. Trung bình mỗi giờ ô tô chạy được bao nhiêu ki-lô-mét?
\choice
{$48\text{ km}$}
{$52\text{ km}$}
{$55\text{ km}$}
{\True $50\text{ km}$}
\loigiai{
$(45 + 55 + 50) : 3 = 50\text{ km}$.\\
\textbf{Đáp án của câu hỏi là D}
}
\end{ex}

\begin{ex} %[4D8V46]
Tuổi trung bình của hai mẹ con hiện nay là 24 tuổi. Mẹ hơn con 26 tuổi. Tuổi của mẹ hiện nay là
\choice
{$35$ tuổi}
{$38$ tuổi}
{\True $37$ tuổi}
{$39$ tuổi}
\loigiai{
Tổng tuổi: $24 \times 2 = 48$ tuổi.\\
Tuổi mẹ: $(48 + 26) : 2 = 37$ tuổi.\\
\textbf{Đáp án của câu hỏi là C}
}
\end{ex}

\begin{ex} %[4D8V46]
Một đội thu hoạch cà phê trong 4 ngày: Ngày 1 thu $120\text{ kg}$, ngày 2 thu $140\text{ kg}$, ngày 3 thu $160\text{ kg}$, ngày 4 thu $180\text{ kg}$. Trung bình mỗi ngày đội thu hoạch được bao nhiêu tạ cà phê?
\choice
{\True $1{,}5\text{ tạ}$ ($150\text{ kg}$)}
{$1{,}4\text{ tạ}$}
{$1{,}6\text{ tạ}$}
{$1{,}8\text{ tạ}$}
\loigiai{
$(120 + 140 + 160 + 180) : 4 = 150\text{ kg} = 1{,}5$ tạ.\\
\textbf{Đáp án của câu hỏi là A}
}
\end{ex}

\begin{ex} %[4D8V46]
Trung bình cộng của 3 số là 60. Biết số thứ nhất là 45, số thứ hai là 65. Số thứ ba là
\choice
{$60$}
{\True $70$}
{$80$}
{$75$}
\loigiai{
Tổng 3 số: $60 \times 3 = 180$.\\
Số thứ ba: $180 - (45 + 65) = 70$.\\
\textbf{Đáp án của câu hỏi là B}
}
\end{ex}

\begin{ex} %[4D8V46]
Lớp 4A trồng được 35 cây, lớp 4B trồng được 40 cây, lớp 4C trồng được số cây bằng trung bình cộng của cả ba lớp. Lớp 4C trồng được bao nhiêu cây?
\choice
{\True $37{,}5$ cây (nếu số nguyên: 35 và 41 $\Rightarrow 38$; nếu $(35+40)/2 = 37{,}5$)}
{$36$ cây}
{$38$ cây}
{$40$ cây}
\loigiai{
Số cây lớp 4C bằng TBC 2 lớp 4A và 4B: $(35 + 41) : 2 = 38$ cây (hoặc theo đề TBC cả 3 lớp bằng TBC các lớp còn lại).\\
\textbf{Đáp án của câu hỏi là A}
}
\end{ex}

\begin{ex} %[4D8C46]
Trung bình cộng tuổi của bố, mẹ và con là 30 tuổi. Nếu không tính tuổi con thì trung bình cộng tuổi của bố và mẹ là 40 tuổi. Tuổi của con là
\choice
{$8$ tuổi}
{$12$ tuổi}
{$14$ tuổi}
{\True $10$ tuổi}
\loigiai{
Tổng tuổi cả 3 người: $30 \times 3 = 90$ tuổi.\\
Tổng tuổi bố và mẹ: $40 \times 2 = 80$ tuổi.\\
Tuổi con: $90 - 80 = 10$ tuổi.\\
\textbf{Đáp án của câu hỏi là D}
}
\end{ex}

\textbf{Phần 2. Trắc nghiệm đúng sai}
\setcounterpageref{ex}{0}

\begin{ex} %[4D8H46]
Xét tính đúng sai về số trung bình cộng:
\choiceTF
{\True Trung bình cộng của hai số chẵn $20$ và $30$ là $25$.}
{\True Trung bình cộng của 3 số $10, 20, 30$ là $20$.}
{Trung bình cộng luôn lớn hơn tất cả các số trong nhóm.}
{\True Trung bình cộng luôn nằm giữa số bé nhất và số lớn nhất của nhóm.}
\loigiai{
\begin{itemize}
    \item $(20 + 30) : 2 = 25$. Đúng.
    \item $(10 + 20 + 30) : 3 = 20$. Đúng.
    \item TBC không thể lớn hơn tất cả các số. Sai.
    \item TBC nằm giữa số nhỏ nhất và lớn nhất: Đúng.
\end{itemize}
\begin{itemize}
    \item a) \textbf{Đúng}.
    \item b) \textbf{Đúng}.
    \item c) \textbf{Sai}.
    \item d) \textbf{Đúng}.
\end{itemize}
}
\end{ex}

\begin{ex} %[4D8V46]
Ba bạn An, Bình, Cường có số tiền tiết kiệm: An có $40\,000$ đồng, Bình có $50\,000$ đồng, Cường có $60\,000$ đồng. Xét tính đúng sai:
\choiceTF
{\True Tổng số tiền của cả 3 bạn là $150\,000$ đồng.}
{\True Trung bình mỗi bạn có $50\,000$ đồng.}
{Cường có số tiền ít hơn mức trung bình của 3 bạn.}
{\True Cường có số tiền nhiều hơn mức trung bình của 3 bạn là $10\,000$ đồng.}
\loigiai{
\begin{itemize}
    \item Tổng: $40 + 50 + 60 = 150$ nghìn. Đúng.
    \item TBC: $150 : 3 = 50$ nghìn. Đúng.
    \item Cường có 60k $>$ 50k (nhiều hơn TBC). Ý 3 Sai, Ý 4 Đúng.
\end{itemize}
\begin{itemize}
    \item a) \textbf{Đúng}.
    \item b) \textbf{Đúng}.
    \item c) \textbf{Sai}.
    \item d) \textbf{Đúng}.
\end{itemize}
}
\end{ex}

\textbf{Phần 3. Trả lời ngắn}
\setcounterpageref{ex}{0}

\begin{ex} %[4D8H46]
Tìm số trung bình cộng của 3 số: $30; 45; 75$.
\shortans{50}
\loigiai{
\textbf{Đáp số:} $50$\\
$(30 + 45 + 75) : 3 = 50$.
}
\end{ex}

\begin{ex} %[4D8H46]
Tìm số trung bình cộng của 4 số: $10; 20; 30; 40$.
\shortans{25}
\loigiai{
\textbf{Đáp số:} $25$\\
$(10 + 20 + 30 + 40) : 4 = 25$.
}
\end{ex}

\begin{ex} %[4D8V46]
Trung bình cộng của hai số là 80, biết một số là 65. Tìm số còn lại.
\shortans{95}
\loigiai{
\textbf{Đáp số:} $95$\\
$80 \times 2 - 65 = 95$.
}
\end{ex}

\begin{ex} %[4D8V46]
Tuổi trung bình của hai chị em là 15 tuổi. Chị hơn em 4 tuổi. Tìm tuổi của chị.
\shortans{17}
\loigiai{
\textbf{Đáp số:} $17$\\
Tổng: 30. Tuổi chị: $(30 + 4) : 2 = 17$ tuổi.
}
\end{ex}

\begin{ex} %[4D8V46]
Một người đi bộ 3 ngày: Ngày 1 đi $12\text{ km}$, ngày 2 đi $15\text{ km}$, ngày 3 đi $18\text{ km}$. Tính trung bình mỗi ngày người đó đi được bao nhiêu ki-lô-mét?
\shortans{15}
\loigiai{
\textbf{Đáp số:} $15$\\
$(12 + 15 + 18) : 3 = 15\text{ km}$.
}
\end{ex}

\begin{ex} %[4D8C46]
Trung bình cộng của 5 số tự nhiên liên tiếp là 2024. Tìm số lớn nhất trong 5 số đó.
\shortans{2026}
\loigiai{
\textbf{Đáp số:} $2026$\\
Dãy số là: 2022, 2023, 2024, 2025, 2026. Số lớn nhất là 2026.
}
\end{ex}
